% ======================================================================
%  Section 4.2.1 --- From least squares to M-estimation.  Corrected.
%
%  FIVE CHANGES. The first is the only one a reader would call an error;
%  the rest are precision, punctuation and a strengthening.
%
%  1. The objective is minimised over V_c in both (4.1) and (4.2). A sum of
%     squared residuals is not a function of a velocity. The cost is a
%     function of the camera POSE; V_c is the descent step the control law
%     produces from it. Corrected below.
%     %%% If Chapters 1 and 3 use the same min_{V_c} notation, they need the
%     %%% same correction, or this section will contradict them.
%  2. "when the weights are held fixed" -- the scale s-hat must be held
%     fixed too, since it is itself a function of the residuals. Without
%     that the stationarity condition and the weighted normal equations do
%     not coincide.
%  3. Equation (4.3) ends in a comma and the following sentence begins with
%     a capital. It should be a full stop.
%  4. The displaced minimiser is now a measured fact, not an assertion:
%     Section 4.7.2 shows Variant 1 converging 26 pixels away while holding
%     the SMALLEST feature error of the three. A forward reference costs one
%     clause and turns the paragraph from motivation into evidence.
%  5. minimize / minimizer are -ize; the chapter body is -ise elsewhere.
% ======================================================================

\section{Robust Estimation in Direct Visual Servoing}
\label{sec:robust-framework}

\subsection{From least squares to M-estimation}
\label{sec:least-squares}
%%% Section 4.7.2 refers back to this subsection; give it a label.

%%% (1) The minimisation is over the pose. The control law realises a step
%%% on the resulting cost; it is not the variable being optimised.
The control laws of Chapters~\ref{chap:Fundamentals-of-Visual-Servoing}
and~\ref{ch:herpvs} minimise a sum of squared residuals. Writing the Hermite
feature error of Chapter~\ref{ch:herpvs} as
$\mathbf{e}(\mathbf{r})=\mathbf{R}(\mathbf{r})-\mathbf{R}^{*}\in\mathbb{R}^{P}$,
where $\mathbf{r}$ denotes the camera pose and $P$ the dimension of the stacked
feature vector, the underlying objective is
\begin{equation}
    \min_{\mathbf{r}}\;\sum_{j=1}^{P} e_{j}^{2}(\mathbf{r}).
    \label{eq:ls-objective}
\end{equation}
The commanded velocity $\mathbf{V}_{c}$ is the descent step that the control
law produces on this objective, not a variable of it.

The quadratic cost weights every measurement by the square of its residual, so
its sensitivity to a measurement grows without bound as that measurement
departs from the model. A block of grossly corrupted measurements, such as
those produced by an occlusion, therefore dominates the objective, and the
minimiser is displaced away from the value that the valid measurements alone
would indicate. That displacement is measured directly in
Section~\ref{sec:ablation-results}, where the unweighted scheme attains the
smallest residual of the three variants while positioning the camera two
orders of magnitude less accurately.
%%% (4) Added. The last sentence is the whole justification for the chapter
%%% and it is now supported by a number rather than left as an expectation.

M-estimation replaces the quadratic by a function $\rho$ that grows more
slowly for large residuals~\cite{key49},
\begin{equation}
    \min_{\mathbf{r}}\;\sum_{j=1}^{P}\rho\!\left(\frac{e_{j}}{\hat{s}}\right),
    \label{eq:m-objective}
\end{equation}
where $\hat{s}$ is a robust estimate of the scale of the residuals, specified
in Section~\ref{sec:tukey}, and required so that $\rho$ is applied to a
dimensionless quantity and the estimator is invariant to the overall magnitude
of the error. Differentiating~\eqref{eq:m-objective} and writing $\psi=\rho'$
for the influence function, the stationarity condition coincides, \emph{when
the weights and the scale are held fixed}, with the normal equations of the
weighted least-squares problem $\min\sum_{j}w_{j}e_{j}^{2}$ with weights
%%% (2) "and the scale" added. s-hat depends on the residuals through the
%%% median absolute deviation (4.7), so freezing the weights alone does not
%%% make the two conditions agree.
\begin{equation}
    w_{j}=\frac{\psi(e_{j}/\hat{s})}{e_{j}/\hat{s}}.
    \label{eq:irls-weight}
\end{equation}
%%% (3) full stop, was a comma.

The weights, however, are themselves functions of the residuals, which depend
in turn on the solution; this circularity is what necessitates an iterative
treatment, and is the basis of iteratively reweighted least
squares~\cite{key50}, in which the weights are computed from the current
residuals and the resulting weighted linear problem is solved, the two steps
alternating. In a visual servoing loop this inner iteration is not performed
to convergence at each control step; instead the weights are recomputed once
per servoing iteration, so that the reweighting and the servoing proceed
together~\cite{key26}.

It is worth noting at the outset that the weight~\eqref{eq:irls-weight}
depends on the magnitude of the residual alone. The estimator therefore
carries no information about \emph{where} in the image a large residual
arises, nor about what the image contains there. This is the opening exploited
in Section~\ref{sec:hermite-weight}.

% ======================================================================
%  NOT CHANGED, AND WHY
%
%  - "grows more slowly for large residuals" is right as a general statement
%    about M-estimators even though Tukey's rho is constant beyond c; the
%    specific behaviour arrives in 4.3.1 and does not need pre-empting here.
%  - The IRLS paragraph is accurate as it stands, including the point that
%    the inner iteration is not run to convergence, which is the correct
%    description of what the implementation does.
%  - \ref{chap:Fundamentals-of-Visual-Servoing} uses a different label
%    convention from ch:herpvs. Invisible in the output; left alone.
%
%  OPTIONAL
%  P is introduced here and used throughout the chapter without a value.
%  Stating it once, in 4.6.2 with the rest of the configuration, would help:
%  P = 264000, being four Hermite scale responses over the 220 x 300
%  interior region left by the ten-pixel filter border.
% ======================================================================
